%
% 公立はこだて未来大学卒業研究中間報告書[全コース対応版]
%
%         ファイル名："sample.tex"
%
\documentclass[11pt]{jarticle}

\usepackage{afuninfosys}
\usepackage{fancyhdr}
\usepackage{url}
\usepackage{float}
\usepackage[dvipdfmx]{graphicx,color}

\author{% 
b10xxxxx 未来太郎\\指導教員 : 函館一郎
}
%\course{Information Systems Course}
%\course{Advanced ICT Course} %% 高度ICTコースの場合はこちらを使用
%\course{Information Design Course} %% 情報デザインコースの場合はこちらを使用
\course{Complex Systems Course} %% 複雑系科学コースの場合はこちらを使用
%\course{Intelligent Systems Course} %% 知能システムコースの場合はこちらを使用
\title{中間報告書の書き方}
\etitle{How to Write Manuscripts for Midterm Report}
\eauthor{Taro MIRAI}
\abstract{和文は300から400文字で記述すること．}
\keywords{北海道, 函館, 亀田中野, 公立はこだて未来大学}
\eabstract{English should be written in 100 to 150 words.}
\ekeywords{Hokkaido, Hakodate, Kamedanakano, FUN}

\begin{document}
\maketitle
%\vspace*{-.5cm}

\section{背景と目的}

このサンプルは情報システムコースにおける中間報告書の様式について説明したものである．必ずしもこの雛形を使う必要はないが，仕上がりイメージはできる限りこの雛形にあわせること．

用紙サイズはA4，向きは縦とし，上下の余白は30mm、左右の余白は25mmとする．本文には明朝体とTimes New Romanを用いる．ただし，タイトルや章節の見出し，図表のキャプションはゴシック体とする．タイトルは14ポイント，氏名と章の見出しは12ポイント，節の見出しは11ポイント，その他は10ポイントとする．また，和文タイトルから英文キーワードまでは1段，本文は2段で構成とし，1段のセクションは42文字×45行，2段のセクションは20文字×45行とする．

なお，章立てはあくまでも参考であり，これに限らない．

\section{関連研究}

中間報告書の文量は4ページとする．学籍番号をファイル名としたPDFファイル１つにまとめた形で作成すること．提出するファイル名はb10xxxxx.pdfとする．

句読点は「，」，「．」とする．「、」，「。」は使用しない．アブストラクトなど英文表記の部分については，スペルチェックプログラムによるチェックをする．

\section{提案する理論}

\subsection{数式}

数式による記述が必要な場合は，式番号を適切に参照しながらまとめること．

\subsection{図・写真}

読者の理解を助けるため，図や表を効果的に利用すること．図のキャプションは

\begin{figure}[H]
  \includegraphics[scale=0.5]{hoge.png}
  \caption{画像のテスト}
  \label{fig1}
\end{figure}

のように，図の下に記す．表のキャプションは

\begin{center}表1　○○○○\end{center}

のように，表の上に記す．

\section{実験と評価}

\section{考察}

\section{結言}

\begin{thebibliography}{99}
\bibitem{marumaru}
	○○△△, システム情報科学会論文誌, 2, 13-19, 2002.
\bibitem{abc}
	A.B.Cdddddd, J. Systems Information Science, 11, 1145-1159, 2001.
\bibitem{batubatu}
	○○××, □□△△, システム情報科学, ☆☆出版, 1999, 20-21.
\bibitem{efghij}
	E.Fggg and H.Ijjj, Electrical Engineering, KKPress, 2003, 281-284.
\end{thebibliography}
\end{document}
%
%
% EOF 
